\documentclass[12pt]{article}

\usepackage{sbc-template}
\usepackage{graphicx,url}
\usepackage[utf8]{inputenc}
\usepackage[brazil]{babel}

\sloppy

\title{Recupera\c{c}\~ao de Informa\c{c}\~ao para \emph{Product Matching} de Insumos SINAPI}

\author{Igor M. Daudt\inst{1}, Luciana Bencke\inst{1}}

\address{Programa de P\'os-Gradua\c{c}\~ao em Computa\c{c}\~ao (PPGC)\\
  Instituto de Inform\'atica -- Universidade Federal do Rio Grande do Sul (UFRGS)\\
  Porto Alegre -- RS -- Brasil
  \email{igor.daudt@inf.ufrgs.br, luciana.bencke@inf.ufrgs.br}
}

\begin{document}

\maketitle

\begin{abstract}
This work evaluates information retrieval methods (BM25, dense E5, hybrid E5, dense BGE-M3, and hybrid BGE-M3) for product matching of SINAPI items, using datasets with distinct lexical profiles: a real ground truth of 852 queries, derived from 948 query-item correlations from a budgeting platform (PainelConstru, real dataset), and a synthetic benchmark of 48 queries, built through systematic technical-notation substitutions. On the real dataset, BM25 attains the highest MAP (0.797), ahead of hybrid BGE-M3 (0.793), with 24 times lower latency. On the synthetic dataset, built for deliberately low lexical overlap (Jaccard 0.403 vs.\ 0.652 on the real dataset), the pattern reverses: hybrid BGE-M3 beats BM25 by $+0.152$ MAP, the largest semantic advantage observed. The advantage of dense retrievers over lexical ones depends on the type of terminological variation and does not fully hold on documents with higher lexical similarity.

\smallskip
\noindent\textbf{Keywords:} information retrieval; product matching; semantic search; hybrid retrieval; SINAPI.
\end{abstract}

\begin{resumo}
Este trabalho avalia m\'etodos de recupera\c{c}\~ao de informa\c{c}\~ao (BM25, E5 denso, E5 \emph{hybrid}, BGE-M3 denso e BGE-M3 \emph{hybrid}) para \emph{product matching} de insumos do SINAPI, usando conjuntos de dados com perfis lexicais distintos: um \emph{ground truth} real de 852 consultas, derivadas de 948 correla\c{c}\~oes consulta-item de uma plataforma de or\c{c}amenta\c{c}\~ao (PainelConstru), e um \emph{benchmark} sint\'etico de 48 consultas, constru\'ido por substitui\c{c}\~oes sistem\'aticas de nota\c{c}\~ao t\'ecnica. No conjunto real, o BM25 obt\'em o maior MAP (0,797), \`a frente do BGE-M3 \emph{hybrid} (0,793), com 24 vezes menos lat\^encia. No conjunto sint\'etico, constru\'ido para ter baixa sobreposi\c{c}\~ao lexical deliberada (Jaccard 0,403 contra 0,652 no conjunto real), o padr\~ao se inverte: o BGE-M3 \emph{hybrid} supera o BM25 em $+0{,}152$ de MAP, a maior vantagem sem\^antica observada. A vantagem de \emph{retrievers} densos sobre o lexical depende do tipo de varia\c{c}\~ao terminol\'ogica e n\~ao se sustenta integralmente em documentos com maior similaridade lexical.

\smallskip
\noindent\textbf{Palavras-chave:} recupera\c{c}\~ao de informa\c{c}\~ao; \emph{product matching}; busca sem\^antica; recupera\c{c}\~ao h\'ibrida; SINAPI.
\end{resumo}


\section{Introdu\c{c}\~ao}

A recupera\c{c}\~ao eficiente de informa\c{c}\~oes em bases t\'ecnicas \'e um desafio recorrente em sistemas de apoio \`a engenharia, compras p\'ublicas e or\c{c}amenta\c{c}\~ao, agravado pela quantidade de itens e varia\c{c}\~oes terminol\'ogicas t\'ipicas do dom\'inio da constru\c{c}\~ao civil -- um problema pr\'oximo ao de \emph{product matching} em cat\'alogos de produtos \cite{peeters:20}. No Brasil, o SINAPI \'e a principal base de refer\^encia de custos para obras p\'ublicas, mas a busca por itens \'e dificultada pela dist\^ancia entre a forma como usu\'arios descrevem produtos e a descri\c{c}\~ao oficial do cat\'alogo.

Este trabalho avalia m\'etodos lexicais, densos e h\'ibridos de recupera\c{c}\~ao de informa\c{c}\~ao sobre o cat\'alogo SINAPI usando os dois conjuntos descritos no Resumo, com perfis lexicais deliberadamente distintos. Comparar os dois conjuntos permite testar a hip\'otese de que \emph{retrievers} densos t\^em maior vantagem sobre o BM25 quanto maior a varia\c{c}\~ao lexical entre consulta e documento -- e verificar se essa vantagem, bem estabelecida para par\'afrases sem\^anticas controladas, se sustenta igualmente sobre vocabul\'ario real de produ\c{c}\~ao. C\'odigo e dados est\~ao publicamente dispon\'iveis.\footnote{\url{https://github.com/igordaudt/Prod_Match_Eramia_26}}

\section{Materiais e M\'etodos}

\textbf{Conjunto real.} O corpus \'e composto por 5.813 insumos do cat\'alogo SINAPI de jan/2026. O \emph{ground truth} \'e formado por 948 correla\c{c}\~oes reais entre produtos cadastrados por usu\'arios e o item SINAPI correspondente, j\'a validadas como parte do uso operacional da ferramenta (geradas por consultas reais na plataforma PainelConstru \cite{painelconstru:26}, \emph{marketplace} de constru\c{c}\~ao). Agrupadas por par consulta-insumo, essas correla\c{c}\~oes formam 852 consultas de avalia\c{c}\~ao, cada uma com um ou mais insumos relevantes.

\textbf{Conjunto sint\'etico controlado (\emph{ppc\_weak}).} Um segundo corpus, de 4.851 itens do cat\'alogo SINAPI, serve de base a um \emph{benchmark} de 48 consultas sobre 16 itens, constru\'ido por fun\c{c}\~oes de rotulagem que substituem sistematicamente a nota\c{c}\~ao t\'ecnica da descri\c{c}\~ao oficial (ex.\ DN em mm~$\rightarrow$~polegadas; remo\c{c}\~ao do termo ``sold\'avel''/``rosc\'avel''), produzindo consultas com Jaccard deliberadamente baixo em rela\c{c}\~ao ao item-alvo. Por usar numera\c{c}\~ao de item independente da do conjunto real, os dois conjuntos s\~ao avaliados separadamente, cada um sobre seu corpus nativo, e comparados no n\'ivel de m\'etricas agregadas.

\textbf{M\'etodos.} Foram avaliados cinco m\'etodos, implementados localmente: BM25 \cite{robertson:09}; E5 \emph{multilingual} denso \cite{wang:22} (\texttt{multilingual-e5-small}); E5 \emph{hybrid} (E5~+~BM25 via \emph{Reciprocal Rank Fusion} \cite{cormack:09}); BGE-M3 \cite{chen:24} (\texttt{bge-m3}) denso e \emph{hybrid} (BGE-M3~+~BM25 via RRF). Cada consulta retornou at\'e 100 candidatos, em CPU, sem GPU. As m\'etricas (MAP e NDCG@10/@100) s\~ao calculadas por consulta e macro-agregadas, invariantes \`a escala absoluta dos scores. Calculou-se tamb\'em a similaridade de Jaccard entre cada consulta e a descri\c{c}\~ao do item-alvo, nos dois conjuntos.

\section{Resultados}

A Tabela~\ref{tab:resultados} apresenta os resultados sobre as 852 consultas do conjunto real (\emph{limit}~=~100). O BM25 obt\'em o maior MAP (0,7967), \`a frente do BGE-M3 hybrid (0,7932) -- diferen\c{c}a de apenas 0,0035 -- com lat\^encia 24 vezes menor (15,9~ms vs.\ 378,6~ms). O BGE-M3 denso puro (MAP$=0{,}7568$, quarto lugar) mostra que a fus\~ao via RRF com BM25 adiciona $+0{,}0365$ de MAP.

\begin{table}[ht]
\centering
\small
\caption{Retrieval sobre o conjunto real (852 consultas, limit=100)}
\label{tab:resultados}
\begin{tabular}{|l|c|c|c|c|}
\hline
\textbf{M\'etodo} & \textbf{MAP} & \textbf{R@10} & \textbf{R@100} & \textbf{Tempo (ms)} \\
\hline
\textbf{BM25} & \textbf{0,7967} & 0,9726 & 0,9988 & \textbf{15,9} \\
\hline
BGE-M3 hybrid & 0,7932 & \textbf{0,9767} & \textbf{1,0000} & 378,6 \\
\hline
E5 hybrid & 0,7638 & 0,9705 & 0,9977 & 48,7 \\
\hline
BGE-M3 (denso) & 0,7568 & 0,9624 & 0,9994 & 297,1 \\
\hline
E5 (denso) & 0,7116 & 0,9417 & 0,9930 & 283,3 \\
\hline
\end{tabular}
\end{table}

A Tabela~\ref{tab:jaccard} confronta esse resultado com o do \emph{ppc\_weak} e com quartis de Jaccard dentro do conjunto real. No \emph{ppc\_weak}, o BGE-M3 hybrid supera o BM25 em $+0{,}152$ de MAP -- a maior vantagem sem\^antica do estudo; no conjunto real, essa vantagem desaparece (BM25 $0{,}797$ vs.\ BGE-M3 hybrid $0{,}793$). Dentro do conjunto real, por\'em, a rela\c{c}\~ao entre Jaccard e vantagem sem\^antica \'e fraca e n\~ao mon\'otona (Spearman $\approx-0{,}07$; o quartil Q3, n\~ao o de menor Jaccard, tem o maior MAP) -- Jaccard baixo em dados reais tem causas heterog\^eneas, nem sempre resol\'uveis por \emph{embeddings}.

\begin{table}[ht]
\centering
\small
\caption{Jaccard m\'edio $\times$ MAP por m\'etodo, nos dois conjuntos}
\label{tab:jaccard}
\begin{tabular}{|l|c|c|c|c|}
\hline
\textbf{Conjunto} & \textbf{N} & \textbf{Jaccard} & \textbf{BM25} & \textbf{BGE-M3 h.} \\
\hline
\textbf{Sint\'etico (ppc\_weak)} & 48 & \textbf{0,403} & 0,559 & \textbf{0,711} \\
\hline
Real -- Q1 (menor Jaccard) & 218 & 0,285 & 0,709 & 0,714 \\
\hline
Real -- Q2 & 222 & 0,670 & 0,776 & 0,782 \\
\hline
Real -- Q3 & 202 & 0,793 & \textbf{0,886} & 0,866 \\
\hline
Real -- Q4 (maior Jaccard) & 210 & 0,866 & 0,823 & 0,818 \\
\hline
Real -- geral & 852 & 0,649 & 0,797 & 0,793 \\
\hline
\end{tabular}
\end{table}

No sint\'etico, o BGE-M3 denso puro obt\'em MAP$=0{,}723$ -- o maior valor entre os cinco m\'etodos ali. A fus\~ao com BM25 tem efeito \emph{oposto} nos dois conjuntos: ajuda no real (visto acima) e atrapalha levemente no sint\'etico ($0{,}723\rightarrow0{,}711$) -- o ganho da fus\~ao depende de o BM25 isoladamente ter algo \'util a contribuir. Uma modelagem de t\'opicos complementar (BERTopic) identificou 90 t\'opicos no corpus real: a vantagem do BM25 concentra-se em fam\'ilias com terminologia t\'ecnica exata, e a do BGE-M3 hybrid, em fam\'ilias com maior varia\c{c}\~ao comercial.

\textbf{Divis\~ao por mediana de Jaccard.} Para testar a hip\'otese de forma mais direta, uniu-se os dois conjuntos numa base combinada de 900 consultas e dividiu-se pela mediana de Jaccard (0,754) em duas metades. A vantagem sem\^antica (MAP BGE-M3 hybrid $-$ MAP BM25) \'e positiva na metade inferior ($+0{,}023$) e negativa na superior ($-0{,}013$) -- \textbf{hip\'otese confirmada}, com efeito modesto (a metade superior \'e quase s\'o consultas reais, 448 de 450). Para verificar se o padr\~ao depende da m\'etrica, recalculou-se tamb\'em o NDCG@10 e @100 por consulta (relev\^ancia bin\'aria, desconto logar\'itmico por posi\c{c}\~ao) e macro-agregou-se por metade (Tabela~\ref{tab:split}) -- mesma l\'ogica do MAP, invariante \`a escala absoluta do score. O NDCG reproduz o mesmo cruzamento nos dois cortes, com amplitude equivalente ($\approx0{,}029$ contra $0{,}036$ do MAP).

\begin{table}[ht]
\centering
\small
\caption{MAP e NDCG por metade de Jaccard (900 consultas combinadas)}
\label{tab:split}
\begin{tabular}{|l|l|c|c|c|}
\hline
\textbf{M\'etrica} & \textbf{Metade} & \textbf{BM25} & \textbf{BGE-M3 h.} & \textbf{Vant. sem.} \\
\hline
MAP & Inferior ($<0{,}754$) & 0,7165 & 0,7396 & $+0{,}023$ \\
\hline
MAP & Superior ($\geq0{,}754$) & 0,8516 & 0,8382 & $-0{,}013$ \\
\hline
NDCG@10 & Inferior & 0,7769 & \textbf{0,7962} & $+0{,}019$ \\
\hline
NDCG@10 & Superior & \textbf{0,8914} & 0,8813 & $-0{,}010$ \\
\hline
NDCG@100 & Inferior & 0,7891 & \textbf{0,8082} & $+0{,}019$ \\
\hline
NDCG@100 & Superior & \textbf{0,8921} & 0,8821 & $-0{,}010$ \\
\hline
\end{tabular}
\end{table}

\section{Discuss\~ao e Conclus\~ao}

O achado central \'e que a vantagem de \emph{retrievers} densos sobre o lexical -- clara no conjunto sint\'etico e confirmada, por MAP e NDCG, pela divis\~ao por mediana de Jaccard -- n\~ao se sustenta integralmente no conjunto real: baixa sobreposi\c{c}\~ao lexical \'e condi\c{c}\~ao necess\'aria, mas n\~ao suficiente, dependendo tamb\'em do tipo de diverg\^encia terminol\'ogica. Na pr\'atica, o BM25 continua um \emph{baseline} forte e barato para \emph{product matching} t\'ecnico; a escolha de \emph{retriever} deve ser validada sobre consultas reais.

Como limita\c{c}\~oes: as diferen\c{c}as de MAP/NDCG e a divis\~ao por mediana carecem de teste de signific\^ancia estat\'istica formal; a metade superior de Jaccard \'e composta quase s\'o por consultas reais (448/450); o \emph{ground truth} \'e de relev\^ancia bin\'aria, n\~ao graduada; e as correla\c{c}\~oes reais n\~ao foram auditadas especificamente para este experimento. Trabalho futuro: \emph{bootstrap} pareado sobre a diferen\c{c}a de vantagem sem\^antica entre metades.

\bibliographystyle{sbc}
\bibliography{references}

\end{document}
